\documentclass{ctexart}
\usepackage{verse}

\newcommand{\attrib}[1]{%
	\nopagebreak{\raggedleft\footnotesize #1\par}}
\renewcommand{\poemtitlefont}{\normalfont\large\itshape\centering}

\begin{document}
	\begin{titlepage}
		\thispagestyle{empty}
		\centering
		
		\vspace*{3.0cm}
		
		{\Large 須以而重集\par}
		\vspace{0.4cm}
		{\Large Juvenilia\par}
		
		\vspace{1.6cm}
		
		{\normalsize 宋捷 Alexandra Chieh Sloan Sung \par}
		\vspace{0.4cm}
		{\normalsize 字：行之\\號：生白無一人\par}
		
		\vspace{1.5cm}
		
		\begin{tabular}{c}
			{\large 收錄期間：2019--2026\par}\\[0.35cm]
			{\large 寫作年齡：9--16歲\par}
		\end{tabular}
		
		\vfill
		
	\end{titlepage}
	
	\begin{titlepage}
		\thispagestyle{empty}
		\centering
		\vspace*{2.2cm}
		
		\begingroup
		\setlength{\parindent}{0pt}
		\setlength{\parskip}{0pt}
		
		\begin{minipage}{0.88\textwidth}
			{\large \itshape
				不得文思 ， 詞句淺陋 。 因是少時戲語妄言 ，而頗有一番可愛 ; 又因常常即興 ，不免格式隨意 ， 難稱詩詞調韻 ， 便為文集罷 。今整理 , 亦未以寫作時序排版 ，僅以作懷記 。}
		\end{minipage}
		
		\endgroup
		
		\vfill
	\end{titlepage}
	
	\begin{titlepage}
		\thispagestyle{empty}
		
		\begingroup
		\setlength{\parindent}{0pt}
		\setlength{\parskip}{0pt}
		
		\noindent
		{\large \itshape \bfseries 寫的第一首詩}
		
		\vspace{2.0cm} 
		
		\begin{center}
			{\large \itshape 惜春}
			\vspace{0.3cm}
			
			\begin{minipage}{\linewidth}
				\centering
				{\large \itshape
					\begin{tabular}{@{}c@{}}
						梨白桃紅枝頭閙\\
						不負韶華爭光陰\\
						哪知零落何去処\\
						尋遍恰逢翌年春
					\end{tabular}
				}
			\end{minipage}
			
		\end{center}
		
		\endgroup
	\end{titlepage}

	
\end{document}
